\begin{tikzpicture}
\path[use as bounding box] (-.1,-.7) rectangle (21.8,11.8);
\node[heading] at (10.8,11.3) {同一条正常流水线：区域做到 115，最终派生到 105};
\node[note,anchor=west] at (2.4,10.55) {A 从 100 出到 110；B 从 110 接棒，已经出到 115。这里比较同一 L2 高度。};
\foreach \xx/\lab in {4/100--105,8.5/106--108,13/109--110,17.5/111--115}{
  \node[note,align=center] at (\xx,9.75) {高度 \lab};
}
\node[rowlabel,text=blue] at (2.0,8.55) {$R$\\区域 FINAL};
\node[rowlabel,text=green] at (2.0,6.55) {$G$\\全局确认};
\node[rowlabel,text=orange] at (2.0,4.55) {$G'$\\接纳进度};
\node[rowlabel,text=violet] at (2.0,2.55) {$D$\\L1 派生};
\foreach \xx/\nm/\lab in {4/r0/共同前缀,8.5/r1/A 的区域结果,13/r2/A 的末块 110,17.5/r3/B 的后继结果}{
  \node[box,text width=3.25cm,fill=regionA] (\nm) at (\xx,8.55) {\lab};
}
\node[box,text width=3.25cm,fill=regularB] at (17.5,8.55) {B 的后继结果};
\draw[flow] (r0) -- (r1); \draw[flow] (r1) -- (r2); \draw[flow] (r2) -- (r3);
\node[box,fill=settled,kept] (g0) at (4,6.55) {共同前缀};
\node[box,fill=settled,kept] (g1) at (8.5,6.55) {已进入 $G$};
\draw[flow] (g0) -- (g1);
\node[note,anchor=west,text width=7.7cm] at (10.8,6.55) {109--115 尚未进入 $G$。\\Publisher / Espresso 故障会卡住这一步。};
\node[box,fill=regularB,draw=orange] (gp0) at (4,4.55) {批次已接纳：cursor = 105};
\node[note,anchor=west,text width=11cm] at (6.5,4.55) {106--108 还未进入 Inbox。\\Batcher / Certifier / DA 故障会卡住这一步。};
\node[box,fill=native,draw=violet] (d0) at (4,2.55) {已派生：确定到 105};
\node[note,anchor=west,text width=11cm] at (6.5,2.55) {正常时，$D$ 依次得到同样的块。\\超窗后，$D$ 可用 deposit-only 越过或偏离 $G$。};
\node[card,fill=gray!3,text width=17.6cm] at (10.9,.55) {
  \textbf{正常时：$D\preceq G'\preceq G\preceq R$。故障后，这个前缀关系可能失效。}\\
  $G'$ 记录接纳进度，不能代替执行状态；恢复时保留的依据是从已最终化 L1 派生出的 $D$。
};
\end{tikzpicture}
